%========================== ARCHITECTURE TECHNIQUE ==========================
\chapter{Architecture technique}

\section{Vue d'ensemble}

L'application repose sur deux briques indépendantes qui communiquent en
\textbf{HTTP/JSON} : un \textbf{frontend React (Vite)} et un \textbf{backend
FastAPI}. Une base de données \textbf{PostgreSQL} est préparée (modèles et
migrations) mais n'est pas branchée par défaut : le backend sert des données
\emph{mockées} en mémoire afin de rester exécutable sans infrastructure de base
de données. La figure~\ref{fig:archi} schématise cette organisation.

\begin{figure}[H]
\centering
\begin{tikzpicture}[
  node distance=1.1cm and 1.6cm,
  box/.style={rectangle, rounded corners=3pt, draw=sbsnavy, line width=0.8pt,
              fill=sbslight, text width=3.5cm, align=center, minimum height=1.5cm,
              font=\small},
  dbbox/.style={cylinder, shape border rotate=90, aspect=0.25, draw=sbsgray,
                dashed, line width=0.8pt, fill=white, text width=3cm,
                align=center, minimum height=1.8cm, font=\small\itshape},
  lbl/.style={font=\scriptsize\itshape, text=sbsgray},
  arr/.style={-{Latex[length=2.5mm]}, line width=0.9pt, draw=sbspurple},
]
  \node[box] (front) {\textbf{Frontend}\\ React + Vite\\ \scriptsize localhost:5173};
  \node[box, right=of front] (back) {\textbf{Backend}\\ FastAPI\\ \scriptsize 127.0.0.1:8000};
  \node[dbbox, right=of back] (db) {PostgreSQL\\ (préparée,\\ non branchée)};

  \draw[arr] ([yshift=4pt]front.east) -- node[lbl, above]{requêtes REST} ([yshift=4pt]back.west);
  \draw[arr] ([yshift=-4pt]back.west) -- node[lbl, below]{JSON} ([yshift=-4pt]front.east);
  \draw[arr, dashed, draw=sbsgray] (back.east) -- node[lbl, above, align=center]{à venir} (db.west);

  \node[below=0.5cm of front, lbl, text width=3.8cm, align=center]
    {Navigation par rôle,\\ pages et composants};
  \node[below=0.5cm of back, lbl, text width=3.8cm, align=center]
    {24 endpoints REST,\\ Swagger/OpenAPI,\\ données en mémoire};
\end{tikzpicture}
\caption{Architecture générale de la plateforme Drive-to-Store DSP}
\label{fig:archi}
\end{figure}

\section{Architecture du frontend}

Le frontend est une application \textbf{React 19} servie par \textbf{Vite}. Il
s'appuie sur \textbf{React Router} pour la navigation, \textbf{Tailwind CSS}
pour le style, \textbf{Recharts} pour les graphiques et \textbf{Leaflet /
OpenStreetMap} pour la cartographie. Son organisation suit une convention
claire :

\begin{itemize}
  \item \texttt{routes/} : déclaration des routes et des gardes d'accès
        (\texttt{RequireAuth}, \texttt{RequireRole}) ;
  \item \texttt{pages/} : une page par écran (Login, Dashboard, Campaigns,
        CampaignCreate, StoreSelection, DCO, Reporting, AccountManagement) ;
  \item \texttt{components/} : composants réutilisables regroupés par domaine
        (\texttt{layout}, \texttt{ui}, \texttt{campaigns}, \texttt{stores},
        \texttt{reporting}, \texttt{account}, \texttt{charts}) ;
  \item \texttt{data/} : une couche d'accès aux données par module
        (\texttt{campaignApi.js}, \texttt{storesApi.js}, \texttt{dcoApi.js},
        \texttt{accountApi.js}, \texttt{reportingData.js}, \texttt{mockData.js}) ;
  \item \texttt{lib/} : utilitaires transverses, dont le client HTTP
        (\texttt{api.js}).
\end{itemize}

Une règle de conception importante : les pages et composants ne font
\textbf{jamais} d'appel réseau directement ; ils passent par la couche
\texttt{data/}, ce qui isole la logique d'accès aux données et facilitera son
évolution.

\section{Architecture du backend}

Le backend est construit avec \textbf{FastAPI} et \textbf{Pydantic v2}. Il est
organisé en couches :

\begin{itemize}
  \item \texttt{routes/} : un routeur par ressource (health, users,
        advertisers, account\_settings, stores, store\_import, campaigns,
        campaign\_creation, dco, statistics) ;
  \item \texttt{schemas/} : schémas Pydantic distinguant la lecture
        (\texttt{*Read}) de l'écriture (\texttt{*Create} / \texttt{*Update}) ;
  \item \texttt{services/} : logique métier et accès aux données, dont
        \texttt{mock\_data.py} (source unique de vérité en mémoire) et
        \texttt{draft\_store.py} (persistance réelle des brouillons) ;
  \item \texttt{models/} : modèles \textbf{SQLAlchemy} correspondant au schéma
        PostgreSQL cible ;
  \item \texttt{core/} : configuration (\texttt{config.py}) et énumérations
        partagées (\texttt{enums.py}).
\end{itemize}

Les routes ne contiennent aucune logique métier : elles appellent un service,
interceptent ses exceptions dédiées (par exemple \texttt{UserNotFoundError},
\texttt{DuplicateEmailError}) et les traduisent en réponses HTTP explicites
(404, 409, 422).

\section{Communication et protection des routes}

Le frontend appelle l'API via un client \texttt{fetch} centralisé, dont l'URL
de base est configurable (\texttt{VITE\_API\_URL}, par défaut
\path{http://localhost:8000/api}). Le backend autorise, via une politique
\textbf{CORS}, les origines des serveurs de développement Vite. Côté frontend,
la \textbf{protection des routes} est assurée par deux gardes : l'une vérifie
qu'une session existe, l'autre que le rôle courant est autorisé pour la route
demandée ; à défaut, l'utilisateur est redirigé (vers la connexion ou le
tableau de bord).

\section{Organisation du dépôt GitHub}

Le dépôt est structuré en trois grands ensembles : \texttt{frontend/}
(application React), \texttt{backend/} (API FastAPI) et \texttt{database/}
(migrations SQL), complétés par un dossier \texttt{docs/} regroupant toute la
documentation. Le versionnement par branches et \emph{pull requests}
(figure~\ref{fig:github}) matérialise chaque incrément de développement.
